\documentclass[12pt,a4paper]{article}

\usepackage[T1]{fontenc}
\usepackage[utf8]{inputenc}
\usepackage[brazil]{babel}
\usepackage{lmodern}
\usepackage{microtype}
\usepackage{geometry}
\usepackage{setspace}
\usepackage{parskip}
\usepackage{enumitem}
\usepackage{needspace}
\usepackage[hidelinks]{hyperref}
\usepackage{xurl}

\geometry{left=2.5cm,right=2.5cm,top=2.5cm,bottom=2.5cm}
\setstretch{1.12}
\setlist[enumerate]{leftmargin=*,itemsep=0.45em,topsep=0.35em}
\setlist[itemize]{leftmargin=*,itemsep=0.35em,topsep=0.3em}
\setlength\emergencystretch{4em}
\Urlmuskip=0mu plus 2mu

\newcommand{\Section}[1]{\needspace{12\baselineskip}\section{#1}}
\newcommand{\SubsectionKeep}[1]{\needspace{10\baselineskip}\subsection{#1}}

\newcommand{\Link}[1]{\href{#1}{\nolinkurl{#1}}}

\newenvironment{NoBreakEnumerate}
{\begin{samepage}\begin{enumerate}}
{\end{enumerate}\end{samepage}}

\newenvironment{NoBreakItemize}
{\begin{samepage}\begin{itemize}}
{\end{itemize}\end{samepage}}

\title{Mapeamento de Análogos ao Terranimo no Brasil\\
\large Foco em compactação por tráfego}
\author{}
\date{Data de referência: 24/02/2026}

\begin{document}

\maketitle
\vspace{-1.2em}

\Section{Objetivo}
Consolidar soluções, documentações e publicações que podem servir de base para um protótipo nacional de avaliação de risco de compactação do solo por tráfego de máquinas, no espírito do Terranimo.

\Section{Síntese executiva}
\begin{NoBreakEnumerate}
  \item Os dois análogos brasileiros mais relevantes encontrados até o momento são: \textbf{PredComp} e \textbf{COMPSOHMA}.
  \item ESALQ e EMBRAPA têm iniciativas digitais importantes em solo, mas não foi identificado publicamente um simulador com o mesmo escopo do Terranimo (carga por roda + pressão de pneu + risco em profundidade por camada).
  \item Recomendação imediata: usar \textbf{PredComp/soilphysics} como benchmark técnico e \textbf{COMPSOHMA/Terranimo Light} como benchmark de interface e fluxo de decisão.
\end{NoBreakEnumerate}

\Section{Análogos diretos (Brasil)}

\SubsectionKeep{PredComp (web app / Shiny)}
\Link{https://appsoilphysics.shinyapps.io/PredComp/}

Relevância:
\begin{NoBreakItemize}
  \item Análogo nacional mais próximo em lógica de modelagem de compactação por tráfego.
  \item Baseado no ecossistema \texttt{soilphysics}, com funções para tensão transmitida ao solo e avaliação de risco.
\end{NoBreakItemize}

\SubsectionKeep{COMPSOHMA (web app / Shiny, beta)}
\Link{https://testbeta.shinyapps.io/compsohma/}

Relevância:
\begin{NoBreakItemize}
  \item Ferramenta brasileira com foco prático de apoio à decisão para risco de compactação.
\end{NoBreakItemize}

\Section{Documentação técnica relevante}
\begin{NoBreakEnumerate}
  \item Índice do pacote \texttt{soilphysics} (CRAN):\\
  \Link{https://search.r-project.org/CRAN/refmans/soilphysics/html/00Index.html}

  \item Função \texttt{stressTraffic} (base SoilFlex):\\
  \Link{https://search.r-project.org/CRAN/refmans/soilphysics/html/stressTraffic.html}

  \item Página do projeto \texttt{soilphysics}:\\
  \Link{https://arsilva87.github.io/soilphysics/}

  \item Código-fonte da app PredComp no pacote:\\
  \Link{https://rdrr.io/cran/soilphysics/src/R/PredComp.R}

  \item Página do pacote no CRAN:\\
  \Link{https://cran.r-project.org/web/packages/soilphysics/index.html}

  \item Registro AGRIS associado ao software/artigo:\\
  \Link{https://agris.fao.org/search/en/providers/122535/records/65de3c294c5aef494fdb1c51}
\end{NoBreakEnumerate}

\Section{Artigos e produção científica (aplicação no Brasil)}
\begin{NoBreakEnumerate}
  \item de Lima et al., 2021 (software \texttt{soilphysics}):\\
  \Link{https://doi.org/10.1016/j.still.2020.104824}

  \item Tese USP, 2017 (sistema de predição para solos brasileiros):\\
  \Link{https://www.teses.usp.br/teses/disponiveis/11/11140/tde-09082017-151718/pt-br.php}

  \item Lozano et al., 2013 (SoilFlex em cana no Brasil):\\
  \Link{https://doi.org/10.1016/j.still.2013.01.010}

  \item Soil \& Tillage Research, 2024 (risco de compactação em cana; contexto Terranimo/SoilFlex/TASC):\\
  \Link{https://doi.org/10.1016/j.still.2024.106206}

  \item Science of the Total Environment, 2019 (TASC em sistemas de cana no Brasil):\\
  \Link{https://doi.org/10.1016/j.scitotenv.2019.05.009}
\end{NoBreakEnumerate}

\Section{Evidência institucional brasileira (COMPSOHMA)}
\begin{NoBreakEnumerate}
  \item CONBEA 2023 (PredComp e COMPSOHMA como ferramentas nacionais):\\
  \Link{https://conbea.org.br/anais/publicacoes/conbea-2023/anais-2023/engenharia-de-agua-e-solo-eas-4/3708-ferramentas-computacionais-made-in-brazil-para-predicao-da-compactacao-do-solo-induzida-pelo-trafego-agricola/file}

  \item Registro de PI (BR512023000761-7):\\
  \Link{https://bv.fapesp.br/en/papi-nuplitec/instituicao/1521/}

  \item Página associada ao ecossistema da ferramenta:\\
  \Link{https://renatoagro.wixsite.com/soilphysics}
\end{NoBreakEnumerate}

\Section{Análogos brasileiros com escopo diferente}
\begin{NoBreakEnumerate}
  \item \textbf{ESALQ} (exemplo de plataforma digital de solo, escopo diferente):\\
  \Link{https://pt.esalq.usp.br/banco-de-noticias/grupo-da-esalq-cria-plataforma-online-para-an\%C3\%A1lise-de-solo}

  \item \textbf{EMBRAPA} (conteúdo e apps sobre solos/manejo, sem equivalência pública direta ao Terranimo no mesmo formato):
  \begin{NoBreakItemize}
    \item \Link{https://www.embrapa.br/busca-de-noticias/-/noticia/47652221/produtor-obtera-classificacao-do-solo-da-sua-propriedade-no-celular}
    \item \Link{https://www.embrapa.br/agencia-de-informacao-tecnologica/cultivos/soja/producao/manejo-do-solo/compactacao-do-solo}
    \item \Link{https://www.embrapa.br/florestas/aplicativos-e-softwares/aplicativos}
  \end{NoBreakItemize}
\end{NoBreakEnumerate}

\Section{Referência internacional (comparação)}
\begin{NoBreakEnumerate}
  \item Página internacional Terranimo por regiões:\\
  \Link{https://www.terranimo.uk/}
\end{NoBreakEnumerate}

Observação: na consulta realizada, não foi identificada versão Brasil explicitamente listada nessa página.

\Section{Conclusão prática para o projeto}
\begin{NoBreakEnumerate}
  \item Benchmark técnico imediato: \textbf{PredComp/soilphysics}.
  \item Benchmark de produto/interface: \textbf{COMPSOHMA + Terranimo Light}.
  \item Próxima etapa: tropicalizar o modelo com parâmetros de natureza e estado dos solos brasileiros (granulometria, Atterberg, índice L, índice de vazios, grau de saturação, histórico de umedecimento/secagem).
\end{NoBreakEnumerate}

\end{document}
